% appendix/questions/concurrentparallel.tex
% mainfile: ../../perfbook.tex
% SPDX-License-Identifier: CC-BY-SA-3.0

\section{"并发"和"并行"有什么区别？}
\label{sec:app:questions:What is the Difference Between ``Concurrent'' and ``Parallel''?}
%
\epigraph{如果你不能制造差别，那就制造区别！}
	 {Unknown}

从经典计算的角度来看，"\IX{concurrent}"和"\IX{parallel}"
显然是同义词。
然而，这并没有阻止许多人在两者之间做出区分，
而且事实证明，这些区分可以从几个不同的角度来理解。

第一种视角将"并行"视为"数据并行"的缩写，
而将"并发"视为几乎所有其他情况。
从这个角度来看，在并行计算中，
整体问题的每个分区可以完全独立地进行，
不需要与其他分区通信。
在这种情况下，分区之间几乎不需要协调。
相比之下，并发计算可能有紧密的相互依赖，
表现为有竞争的锁、事务或其他同步机制。

\QuickQuiz{
	假设程序的一部分仅使用RCU读端原语作为其唯一的同步机制。
	这是并行还是并发？
}\QuickQuizAnswer{
	两者都是。
}\QuickQuizEnd

这当然引出了这种区分为什么重要的问题，
这就把我们带到了第二种视角，即底层调度器的视角。
调度器的复杂性和能力范围很广，
作为一个粗略的经验法则，一组并行进程之间的通信越紧密和不规则，
对调度器所要求的精巧程度就越高。
因此，并行计算避免相互依赖意味着并行计算程序可以在
能力最弱的调度器上良好运行。
事实上，一个纯数据并行程序在被任意细分和交错到
单处理器上后仍然可以成功运行。\footnote{
	是的，这确实意味着数据并行程序最适合串行执行。
	你为什么这么问？}
相比之下，并发计算程序可能需要调度器具有极高的精巧性。

人们可以争辩说，我们应该简单地要求调度器具有合理的能力水平，
这样我们就可以忽略并行和并发之间的任何区别。
虽然这通常是一个好策略，
但在一些重要的情况下，\IX{efficiency}、
\IX{performance}和\IX{scalability}方面的考虑
严重限制了调度器能够合理提供的能力水平。
一个重要的例子是当调度器在硬件中实现时，
正如在SIMD单元或GPGPU中经常出现的那样。
另一个例子是工作单元非常短的工作负载，
即使是基于软件的调度器也必须在精巧性和效率之间做出艰难的选择。

现在，第二种视角可以被认为是使工作负载与可用的调度器匹配，
并行工作负载可以使用简单的调度器，
而并发工作负载需要复杂的调度器。

不幸的是，这个视角并不总是与第一种视角提出的基于依赖关系的区分一致。
例如，一个高度相互依赖的基于锁的工作负载，
每个CPU一个线程，可以用一个简单的调度器就能应付，
因为不需要做任何调度决策。
事实上，这种类型的某些工作负载甚至可以在串行机器上一个接一个地运行。
因此，这样的工作负载会被第一种视角标记为"并发"，
而被许多采取第二种视角的人标记为"并行"。

\QuickQuiz{
	在第二种（基于调度器的）视角的哪个部分，
	基于锁的单线程每CPU工作负载会被视为"并发"？
}\QuickQuizAnswer{
	那些希望任意细分和交错工作负载的人会这样认为。
	当然，任意细分可能最终将锁的获取与对应的锁释放分开，
	这将阻止任何其他线程获取该锁。
	如果这些锁是纯自旋锁，这甚至可能导致死锁。
}\QuickQuizEnd

第三种视角认为并发是源代码的逻辑表现，
而并行是在实际硬件上运行该代码的物理表现。
但那么，对于根据可用CPU数量生成可变数量线程的用户应用程序，
我们该如何看待？
这些应用程序在单个CPU上运行时是并发的，
而在多个CPU上运行时是并行的吗？
如果这些应用程序即使在多个CPU上仍然是并发的而非并行的，
那如果它们使用\co{sched_setaffinity()}等工具
将每个线程绑定到指定的CPU上呢？
那么，这些应用程序是硬编码CPU编号，
还是根据可用CPU计算CPU编号，会有区别吗？
如果这些应用程序仍然只是并发的，
那如果它们在少于（比如说）四个CPU可用时拒绝运行呢？
如果瓶颈是CPU以外的其他东西呢，
正如\cref{sec:datastruct:Hash-Table Performance}中所讨论的？
当这些问题及其变体在生产环境中运行此类应用程序的经验下
被仔细考虑时，请注意本书中的可扩展性图表
完全证明了书名中"并行"一词的使用。

简言之，这种视角可能具有易于教学的优点，
但从实际角度来看似乎有所欠缺。

这完全没问题。
毕竟，任何人类视角在客观宇宙面前都不具有分量。

这种分类失败并不一定意味着这些视角毫无用处，
而是说当你考虑是否尝试应用它们时，
应该持有适当怀疑的态度。
一如既往，在适用的地方使用给定的视角，否则忽略它。

事实上，除了并行、并发、map-reduce、基于任务等之外，
可能还会出现新的视角。
有些会经受住时间的考验，有些不会。
祝你好运，猜猜哪些能经受考验，哪些不能！
